\section{Ugello}
Come nel caso subsonico, è stato inizialmente assunto di essere in un caso di espansione ottima ma in seguito alle prime analisi, a causa di una quota di funzionamento maggiore, l'area di efflusso dell'ugello è risultata troppo alta; è stato quindi considerato un ugello che permette di avere un flusso sottoespanso in uscita.
Il valore di pressione in uscita è stato quindi fissato a \(4\) volte la pressione ambiente a \(h=25 \text{ km}\).
Il resto del dimensionamento è stato poi eseguito con le stesse modalità del caso subsonico.

\begin{table}[H]
    \centering
    \caption{Valori ottenuti per l'ugello}
    \renewcommand{\arraystretch}{1}
    \fontsize{11pt}{11pt}\selectfont

    \begin{tabular}{cllc}
         \toprule
         \textbf{Grandezza} & \textbf{Simbolo} & \textbf{Valore} & \textbf{Unità di misura} \\
         \midrule
         Pressione totale in uscita & $P_{05}$ & $114$ & \text{kPa} \\
         Temperatura totale in uscita & $T_{05}$ & $2100$ & \text{K} \\
         Mach di efflusso & $M_{05}$ & $2.23$ & - \\
         Velocità di efflusso & $V_{e}$ & $1480$ & \text{m/s} \\
         Area di gola & $A^\star$ & $0.891$ & $\text{m}^2$ \\
         Area di efflusso & $A_{e}$ & $1.93$ & $\text{m}^2$ \\
         \bottomrule
    \end{tabular}
\label{tab:ugello_ramjet}
\end{table}
Come detto in precedenza la variazione delle aree di gola e di efflusso è resa possibile grazie ad un sistema di superfici mobili.
%inserire disegno